\section{Multi-Omics Profiling in \gls{lusc} Research: Platforms and Data Types}

%TODO: i have a quick table but i feel like it misses the mark completely; by that, i mean - it doesnt cover how the data was actually generated; like kind of a background, what methods, what sample groups, what technologies, the tcga lusc dataset has metadata.. im talking about that, what demographics, details like that..

\subsection{Overview of TCGA Omics Platforms}

The TCGA project and the UCSC Xena platform have made matched multi-omics cancer data broadly accessible for computational analysis \parencite{Goldman2020}. In \gls{lusc}, five omics modalities are available: gene expression, \gls{dna} methylation, \gls{mirna} expression, copy number variation, and RPPA. Each measures a distinct molecular layer with different biological content, dynamic range, and analytical properties.

%TODO: does this table need references? maybe in caption or rows or sth..

\begin{table}[ht]
\small
\caption[Omics modalities used for LUSC subtype classification]{Omics
modalities used in this thesis and their key properties.}
\centering
\begin{tabular}{p{2.5cm}p{3.0cm}p{3.0cm}p{4.2cm}}
\toprule
\textbf{Modality} & \textbf{Measured layer} & \textbf{Feature space}
    & \textbf{Interpretive value} \\
\midrule
Gene expression & mRNA transcript abundance
    & $\sim$17{,}000--20{,}000 genes after filtering
    & Directly reflects transcriptional programs defining subtype labels. \\[3pt]
\gls{dna} methylation & CpG methylation beta values
    & $\sim$450{,}000 probes; 2{,}000 retained
    & Epigenetic regulation; context-dependent interpretation. \\[3pt]
\gls{mirna} expression & Mature small-RNA abundance
    & $\sim$200 miRNAs after filtering
    & Post-transcriptional regulatory programs and lineage markers. \\[3pt]
Copy number variation & GISTIC2 thresholded calls
    & 2{,}000 most variable genes
    & Structural genomic alteration; discrete ordinal scale. \\[3pt]
RPPA & Protein/phosphoprotein abundance
    & 238 antibodies post-QC
    & Functional signalling states; closest to phenotype. \\
\bottomrule
\end{tabular}
\label{tab:omics_layers_lit_review}
\end{table}


Gene expression data quantify transcript abundance and are the most directly relevant modality for recovering transcriptomically defined subtypes \parencite{Wilkerson2010}. RSEM-normalised RNA-seq values provide gene-level estimates compatible with downstream supervised learning after expression-level filtering and variance-based feature selection \parencite{Li2013}. \gls{dna} methylation arrays measure beta values at CpG sites across promoters, gene bodies, CpG islands, shores, shelves, and intergenic regions \parencite{Bibikova2011}; interpretation depends critically on probe genomic context, as promoter hypermethylation is commonly associated with transcriptional silencing whereas gene-body methylation may have activating or neutral effects. \gls{mirna} expression profiles are compact (hundreds of measured miRNAs) and may offer a more stable classification space than the full transcriptome \parencite{Wilkerson2010}. GISTIC2-thresholded copy number variation data encode discrete genomic gain and loss events \parencite{Mermel2011}; although \gls{lusc} carries recurrent chromosomal alterations, thresholded CNV may not resolve expression-defined subtypes as effectively as regulatory data layers. RPPA measures protein abundance using validated antibody panels, capturing functional pathway activity closer to the cellular phenotype than mRNA abundance alone \parencite{Tibes2006}.


\subsection{Preprocessing Standards in Omics Studies}

%TODO: looks thin, what are the standard preprocessing pipelines or studies such as mine? current section doesnt address that at all.

Appropriate preprocessing is essential for valid supervised learning on omics data and differs substantially across platforms.
 
For RNA-seq gene expression, RSEM-normalised values are filtered to remove unexpressed genes (median log$_2$ expression $< 1.0$) before variance-based feature selection. Median absolute deviation (MAD) is preferred over standard deviation for omics feature selection because it is robust to outlier samples and monotonically related to
biological variability \parencite{Hoadley2018}. The top 2{,}000 most variable genes by MAD were retained in the present study, a threshold widely used in cancer classification studies \parencite{Hoadley2018, Picard2021}.
 
For methylation, Level~3 beta values undergo probe-level quality control: probes missing in more than 20\% of samples are removed, remaining missing values are imputed using per-probe median, and near-invariant probes (standard deviation $< 0.05$) are excluded \parencite{Bibikova2011}. Probes with rs-prefix (SNP probes) are identified separately because their variation reflects genotype rather than epigenetic state, and their presence among top-ranked features should not be interpreted as methylation biology \parencite{Bibikova2011}.
 
For \gls{mirna} expression, missing values are treated as structural zeros (zero reads detected) rather than technical missingness and are filled accordingly \parencite{Wilkerson2010}. Low-abundance miRNAs (median log$_2$(RPM+1) $< 1.0$) are filtered before MAD-based selection of the 200 most variable mature miRNAs. For copy number variation, discrete GISTIC2 calls are used directly without normalisation, as the ordinal scale encodes categorical genomic states \parencite{Mermel2011}. For RPPA, normalised protein abundance values undergo missing-value filtering and median imputation, but no further feature selection is required given the limited panel size (238 antibodies post-QC) \parencite{Tibes2006}.
 
A critical consideration across all modalities is the prevention of information leakage \parencite{Vabalas2019}. Any data-adaptive step - feature scaling, missing-value imputation from training data, scaler fitting - must be performed exclusively on training folds and applied to validation folds using parameters estimated from training data only.
In this thesis, StandardScaler normalisation was fitted and applied within training folds; variance-based MAD filtering was applied once before cross-validation as an unsupervised preprocessing step, which is a minor limitation acknowledged in the discussion.


\subsection{Modality-Specific Biological Considerations}

Each omics layer provides a different window onto tumour biology. Gene expression is expected to most faithfully encode the expression-defined subtype taxonomy \parencite{Wilkerson2010}. Methylation provides orthogonal epigenetic information but requires promoter-context annotation for mechanistic interpretation, which is a known limitation in purely data-driven analyses \parencite{Bibikova2011}. \gls{mirna} profiles regulate mRNA stability and translation and may capture lineage-state information compactly \parencite{Wilkerson2010}. RPPA captures functional signalling
pathway activity but is limited by the antibody panel breadth \parencite{Tibes2006}. CNV captures structural genomic alterations that may be shared across subtypes and may not preserve expression- defined subtype boundaries \parencite{Hammerman2012}.
 
Multi-omics cancer studies increasingly pursue integration rather than modality-by-modality comparison \parencite{Hoadley2018, Picard2021, Nicora2020, Argelaguet2020}. Integration can improve classification and reveal inter-modality dependencies, but it obscures which platform contributes which biological information. The complementary approach taken here - independent modality
comparison before integration - establishes a principled baseline for assessing each platform's discriminative value.
